\section{京大過去問集}
\subsection{基礎問題}
\begin{prob} \checkbox (H28基礎科目II) \\
$A = \bmat{0&-1\\ 1&1}$, $B = \bmat{0&1\\ 1&0}$で生成される$GL_2(\mb{C})$の部分群$G$について以下の問に答えよ.
\ben
\item $|G|$を求めよ.
\item $|Z(G)|$を求めよ.
\item $G$の位数2の元の個数を求めよ.
\een
\end{prob}
\begin{point}
生成元が二つあたえられているので, とりあえずいろいろ計算して関係式を出してみる. 共役元を求めるといいことがたまにある.  (1)の問題の出し方的に, $A,B$は有限位数であることにも注意.
\end{point}
\begin{ans}
$A^6 = B^2 = I$はすぐ分かる. $BAB^{-1} = \sumib{1&1\\ 0 & -1}\sumib{0&1\\ 1&0} = \sumib{1&1\\ -1&0} = A^{-1} = A^5$である. よって$BA = A^{-1}B$. $G$の元は, $A^6 = B^2 = I$に注意すると
\[A^{k_1} BA^{k_2}B\dots A^{k_l}B^{b},\quad (0\leq k_i\leq 5, b=0,1)\]
の形に書けるが, $BA= A^{-1}B$を用いると必ず$A^{a}B^b$の形に書ける(ちゃんとやりたいなら$l$に関する帰納法を用いる). ただし$0\leq a\leq 5, b=0,1$. よって, この表記法が12通りしかないので$|G| \leq 12$がわかる. 一方で$A$の位数が6なので$|G|$は6の倍数であり, $B$は$A$の冪にならないので$|G|$は6より大きい. よって$|G| = 12$が従う. \\
(2): $|G|=12$より, $G$の元を$A^aB^b$ ($a=0,\dots, 5, b=0,1$)の形に表す方法は一意. 実は, このような表記に関して$b=1$であるものは中心に属さない. つまり, $A^aB\in Z(G)$となることはない.  なぜならば, もし$A^aB\in Z$なら
\[A(A^a B) = A^{a+1}B = (A^a B)A = A^{a-1}B\]
となり表記が一意であることに反するから. 次に $A^a\in Z(G)$であるための必要条件を考えよう. $BA=A^{-1}B$より$ABA = B$なので, $B=ABA = A(ABA)A = A^2BA^2 = \dots = A^nB A^n$である. よって,
\[A^a B = BA^a = A^{-a} (A^a B A^a) = A^{-a} B \]
となり, 表記の一意性から$a\equiv -a\pmod{6}$である. よって$a=0,3$が必要. 逆にこのとき$A^a = \pm I$なので明らかに中心に属す. よって$|Z(G)| = 2$. \\
(3): $A^aB$の形の元は位数2である. 実際, $A^a B \neq I$で
\[(A^aB)^2 = (A^aB A^a)B = B B = I\]
1であるから. よって$A^aB$ ($a=0,1,2,3,4,5$)は位数2. また, $A$は位数6なので, $A^a$の形の元で位数2のものは$A^3$のみである. 以上より$7$個ある.
\end{ans}



\subsection{専門問題}
\begin{prob} \checkbox (R2専門科目)
位数$3p^n$ ($p$は素数) の群は可解であることを示せ.
\end{prob}
\begin{hint}
\chatgpthint{Burnsideの$p^aq^b$定理を適用してください。}
\end{hint}

\begin{prob} \checkbox (H29専門科目)
$G$を有限群とする.
\ben
\item $|\mr{Aut}(G)| = 1$のとき, $G$は自明群であるか, または$\mb{Z}/2\mb{Z}$と同型であることを示せ.
\item $|G|$が奇数で$|\mr{Aut}(G)| = 2$であるような$G$をすべて求めよ.
\een
\end{prob}
\begin{hint}
\chatgpthint{内部自己同型と反転写像が自己同型になる条件を調べてください。}
\end{hint}

\begin{prob} \checkbox (H28専門科目)\\
$G = GL_2(\mb{F}_{p})$とおく.
\ben
\item $G$の元で対角行列と共役でないものの個数を求めよ.
\item $X,Y\in G$の最小多項式が一致すれば$X,Y$は互いに共役であることを示せ.
\item 対角行列を含まない$G$の共役類の個数を求めよ.
\een
\end{prob}
\begin{hint}
\chatgpthint{$2$次行列の有理標準形を最小多項式ごとに分類してください。}
\end{hint}

\begin{prob} \checkbox (H27専門科目)
$G$は非可換群で「$N_1,N_2\subsetneq G$が相異なる$\set{1},G$以外の正規部分群なら, $N_1\not\subset N_2$」という条件を満たすとする.
\ben
\item$N_1,N_2$が相異なる$\set{1},G$以外の正規部分群なら, $G=N_1\times N_2$であることを証明せよ.
\item $G$の自明でない正規部分群の数は高々2個であることを証明せよ.
\een
\end{prob}
\begin{hint}
\chatgpthint{相異なる正規部分群の積と共通部分を調べてください。}
\end{hint}

\begin{prob} (H26専門科目)
$K\subset \mb{C}$を部分体, $p$を素数とする. $\mb{C}$に含まれる任意の有限次拡大$L/K$に対し, $L=K$でなければ$[L:K]$は$p$で割り切れると仮定する. このとき, $\mb{C}$に含まれる任意の有限次拡大$L/K$に対し$[L:K]$は$p$の冪(1を含む)であることを証明せよ.
\end{prob}
\begin{hint}
\chatgpthint{正規閉包のGalois群とSylow $p$部分群の固定体を考えてください。}
\end{hint}
